\chapter{Exact streamed representations}\label{ch:streaming}
\chapterlead{Large factors need not be resident all at once. The identities in this
chapter separate total data size from the working set without changing the
operator. A storage partition is therefore a choice about execution, not a
statistical approximation. Weighted features provide a concrete construction,
but every subsequent identity applies to a supplied real factor.}
\section{Weighted-feature specialization}\label{sec:final-7}

A prominent application constructs the general factor from weighted,
covariate-residualized features.  The sequence kernel association test (SKAT) provides one application
\cite{Wu2011}, but the same formulas apply to any feature matrix with
nonnegative column weights.

Let
\[
G\in\mathbb R^{n\times m},
\qquad
W=\operatorname{diag}(w_1,\ldots,w_m),
\quad w_j\ge 0,
\]
and let \(Q_X\in\mathbb R^{n\times p}\) have orthonormal columns spanning
the null-model covariates.  Define
\[
M_X=I-Q_XQ_X^\top,
\qquad
A=M_XGW^{1/2},
\qquad
K=AA^\top.
\]

The central implementation difficulty is not the algebraic size of the
model but the mismatch between the total factor size \(n\times m\) and the
working memory of an accelerator.  The key point of this book is that
several decompositions act on different mathematical coordinates and therefore
compose:
\[
\boxed{
\text{storage blocks}
\circ
\text{structural phases}
\circ
\text{Krylov/range compression}
\circ
\text{spectral tail packets}.
}
\]

Storage blocking controls residence.  Structural phases encode linkage-disequilibrium (LD) windows,
annotations, frequency bands, or other coordinate groups.  Krylov methods
compress the operator to the rank actually requested.  Spectral packets
compress the final generalized-\(\chi^2\) transform.


\begin{figure}[htbp]
\centering
\includegraphics[width=\linewidth]{figures_production/figure_01_residualization_geometry.pdf}
\caption{Orthogonal residualization separates a feature into a component in the covariate span and a component retained by $M_X$. The diagram is an illustrative coordinate model. The identity behind it is the exact projector decomposition.}
\label{fig:residualization-geometry}
\end{figure}


\section{Variant-block Gram decomposition}\label{sec:final-8}

Partition the variant indices into disjoint blocks
\[
\{1,\ldots,m\}
=
J_1\sqcup\cdots\sqcup J_B.
\]
Let
\[
G_b=G_{[:,J_b]},
\qquad
W_b=W_{J_b,J_b},
\qquad
A_b=M_XG_bW_b^{1/2}.
\]
Then
\[
A=[A_1\ \cdots\ A_B].
\]

\begin{theorem}[Additive block Gram identity]\label{claim:final-8-1}
\[
\boxed{
K
=
AA^\top
=
\sum_{b=1}^B A_bA_b^\top
=
M_X
\left(
\sum_{b=1}^B G_bW_bG_b^\top
\right)
M_X.
}
\]
\end{theorem}

\begin{proof}
Horizontal concatenation gives
\[
[A_1\ \cdots\ A_B]
\begin{bmatrix}
A_1^\top\\
\vdots\\
A_B^\top
\end{bmatrix}
=
\sum_bA_bA_b^\top.
\]
The final form follows from
\(A_b=M_XG_bW_b^{1/2}\) and
\(M_X^\top=M_X=M_X^2\).
\end{proof}

\begin{remark}\label{claim:final-8-2}
The theorem is exact for every partition.  Block boundaries are execution
parameters, not statistical parameters.
\end{remark}


\begin{figure}[htbp]
\centering
\includegraphics[width=\linewidth]{figures_production/figure_02_exact_block_kernel_assembly.pdf}
\caption{Each resident column block contributes a complete sample-space Gram summand $A_bA_b^{\mathsf T}$. Adding the summands recovers $AA^{\mathsf T}$ exactly. Storage partitioning changes residence and evaluation order, not the represented operator.}
\label{fig:block-kernel-assembly}
\end{figure}


\section{Exact streamed kernel actions}\label{sec:final-9}

\begin{theorem}[One-pass sample-space streaming action]\label{claim:final-9-1}
For every block vector \(V\in\mathbb R^{n\times s}\),
\[
\boxed{
KV
=
M_X
\sum_{b=1}^B
G_bW_bG_b^\top
M_XV.
}
\]
\end{theorem}

\begin{proof}
Apply the additive block Gram identity to \(KV\), then associate each block
product as
\[
G_bW_b(G_b^\top M_XV).
\]
\end{proof}

\begin{algorithm}[Streamed \(KV\)]\label{claim:final-9-2}
\begin{enumerate}[label=\arabic*.]
\item Compute \(U=M_XV\).
\item Set \(Y=0\).
\item For each resident block \(b\), compute
\[
T_b=G_b^\top U,
\qquad
Y\leftarrow Y+G_b(W_bT_b).
\]
\item Return \(M_XY\).
\end{enumerate}
\end{algorithm}

\begin{corollary}[Working memory independent of total variant count]\label{claim:final-9-3}
If at most \(c\) variants are resident and the block-vector width is \(s\),
then the algebraic working set is
\[
\boxed{
O(nc+ns+cs+np),
}
\]
independent of total variant count \(m\).
\end{corollary}


\section{Streamed factor actions and Golub--Kahan}\label{sec:final-10}

Let \(V\in\mathbb R^{m\times s}\) be partitioned conformally:
\[
V=
\begin{bmatrix}
V_1\\
\vdots\\
V_B
\end{bmatrix}.
\]

\begin{theorem}[Streamed \(AV\)]\label{claim:final-10-1}
\[
\boxed{
AV
=
M_X
\sum_{b=1}^B
G_bW_b^{1/2}V_b.
}
\]
\end{theorem}

\begin{theorem}[Streamed \(A^\top U\)]\label{claim:final-10-2}
For \(U\in\mathbb R^{n\times s}\),
\[
\boxed{
(A^\top U)_b
=
W_b^{1/2}G_b^\top M_XU.
}
\]
\end{theorem}

\begin{proof}
Both identities follow by substituting
\(A=M_XGW^{1/2}\) and partitioning the variant coordinate.
\end{proof}

\begin{corollary}[Streaming Golub--Kahan equivalence]\label{claim:final-10-3}
Any exact-arithmetic algorithm whose access to \(A\) occurs only through
\(AV\) and \(A^\top U\) produces the same iterates when those products are
implemented by the preceding streamed formulas.
\end{corollary}

This applies directly to block Golub--Kahan bidiagonalization, matrix-free
singular-value iteration, range finding, and restarted variants.  The
factor need not be resident.


\section{Two-pass dual action}\label{sec:final-11}

The variant-space normal operator is
\[
H=A^\top A.
\]

\begin{theorem}[Two-pass streamed dual action]\label{claim:final-11-1}
For every \(V\in\mathbb R^{m\times s}\),
\[
\boxed{
HV=A^\top(AV).
}
\]
The product is evaluated exactly by a streamed pass forming
\(Y=AV\in\mathbb R^{n\times s}\), followed by a second streamed pass
emitting the blocks
\[
(HV)_b=W_b^{1/2}G_b^\top M_XY.
\]
\end{theorem}

This form is particularly natural for variant-space polyphase Krylov
iterations because the right-coordinate blocks can be produced and consumed
sequentially.


\section{Two-dimensional sample/variant tiling}\label{sec:final-12}

A variant block may itself exceed device memory.  Let
\[
H_b=G_bW_b^{1/2}
\]
and partition samples into row tiles \(I_1,\ldots,I_R\).

\begin{theorem}[Two-pass residualized tile construction]\label{claim:final-12-1}
Define
\[
R_b=Q_X^\top H_b.
\]
Then
\[
\boxed{
R_b
=
\sum_r
Q_{X,I_r}^\top H_{b,I_r},
}
\]
and after \(R_b\) is known,
\[
\boxed{
A_{b,I_r}
=
H_{b,I_r}-Q_{X,I_r}R_b.
}
\]
Consequently,
\[
\boxed{
K_{I_r,I_s}
=
\sum_b
A_{b,I_r}A_{b,I_s}^\top.
}
\]
\end{theorem}

\begin{proof}
The first identity is the row-partitioned matrix product
\(Q_X^\top H_b\).  The second is the definition
\(M_XH_b=H_b-Q_XQ_X^\top H_b\).  The last identity is the corresponding
sample-tile expansion of \(AA^\top\).
\end{proof}

With sample tile height \(t\) and resident variant width \(c\), the dominant
working storage, including a resident covariate tile, is
\[
O(tc+pc+tp+t^2).
\]
This count excludes the backing input store, completed output blocks,
allocator overhead, and library workspaces.


\section{Partition refinement and permutation invariance}\label{sec:final-13}

\begin{theorem}[Partition-refinement invariance]\label{claim:final-13-1}
Let \(\mathcal J\) be any partition of the variants and let
\(\mathcal J'\) refine \(\mathcal J\).  Then
\[
\boxed{
\sum_{J\in\mathcal J}A_JA_J^\top
=
\sum_{J'\in\mathcal J'}A_{J'}A_{J'}^\top
=
AA^\top.
}
\]
The analogous streamed operator action is also invariant.
\end{theorem}

\begin{proof}
Each coarse block is a disjoint union of refined blocks.  Apply the additive
Gram theorem within each coarse block and sum.
\end{proof}

\begin{corollary}[Adaptive memory splitting]\label{claim:final-13-2}
A runtime may recursively subdivide any block that exceeds available device
memory without changing the exact operator.
\end{corollary}

\begin{proposition}[Variant permutation invariance]\label{claim:final-13-3}
If the columns of \(G\) and the diagonal entries of \(W\) are permuted by the
same permutation, \(K\) and every streamed sample-space action are unchanged.
\end{proposition}


\begin{figure}[htbp]
\centering
\includegraphics[width=\linewidth]{figures_production/figure_08_partition_invariance.pdf}
\caption{Gram and operator-action discrepancies for five retained partition controls. Horizontal connectors pair the two measurements for each partition. The small floating discrepancies support the implementation check, while the partition identity itself is exact algebra.}
\label{fig:partition-invariance}
\end{figure}


\section{Memory planning law}\label{sec:final-27}

Let \(d\) be the number of bytes per floating scalar, \(c\) the resident
variant count, \(s\) the Krylov width, and \(p\) the covariate rank.
Ignoring library workspace and allocator overhead, a streamed sample-space
action has working memory on the order of
\[
\boxed{
d\,[nc+2ns+cs+np+O(ps+c)].
}
\]

The full factor costs
\[
dnm.
\]

For fixed \(n,p,s,c\),
\[
\boxed{
\frac{M_{\rm stream}}{M_{\rm full}}
=
O\!\left(
\frac{c}{m}
+
\frac{s}{m}
+
\frac{p}{m}
\right),
}
\]
so device memory remains bounded as total variant count grows.

\begin{table}[htbp]
\centering
\caption{Algebraic float64 memory model for \(n=2000\), \(p=4\),
Krylov width \(s=16\), resident variant block \(c=1024\).}
\begin{tabular}{rrrr}
\toprule
\(m\) & full factor (MiB) & streamed \(KV\) (MiB) & tiled Gram (MiB)\\
\midrule
10,000 & 152.6 & 16.308 & 8.555\\
50,000 & 762.9 & 16.308 & 8.555\\
100,000 & 1,525.9 & 16.308 & 8.555\\
500,000 & 7,629.4 & 16.308 & 8.555\\
1,000,000 & 15,258.8 & 16.308 & 8.555\\
\bottomrule
\end{tabular}
\end{table}


\begin{figure}[htbp]
\centering
\includegraphics[width=\linewidth]{figures_production/figure_03_memory_law.pdf}
\caption{Binary64 payload as total factor width increases. The complete factor grows with $m$, while the streamed and tiled payloads remain constant at fixed resident dimensions. These are algebraic sizing data, not measurements of total device allocation.}
\label{fig:memory-law}
\end{figure}



\section{The running example in three storage blocks}\label{sec:worked-stream}
Partition the factor of Section~\ref{sec:worked-law} into consecutive pairs of columns,
\[
A=[A_1\ A_2\ A_3],\qquad A_i\in\mathbb R^{4\times2}.
\]
The second and third blocks are identical, but they are different column positions in the factor. Both must be counted. Their Gram contributions are
\begin{align*}
A_1A_1^\top&=\frac12
\begin{pmatrix}1&0&1&0\\0&1&0&1\\1&0&1&0\\0&1&0&1\end{pmatrix},\\
A_2A_2^\top=A_3A_3^\top&=\frac12
\begin{pmatrix}1&0&-1&0\\0&1&0&-1\\-1&0&1&0\\0&-1&0&1\end{pmatrix}.
\end{align*}
Adding all three gives the matrix $K$ already derived. Merging the last four columns into one storage block gives the same result. Discarding the repeated third block does not: that changes its coefficient contribution and hence the probability law.

For the declared vector $v=(1,2,-1,0)^\top$, the action can be evaluated without forming any block Gram matrix:
\[
\sum_{j=1}^3 A_j(A_j^\top v)=Kv=(2,3,-2,-1)^\top.
\]
Every operation in this identity is rational. The companion calculation checks it symbolically, along with the change from three blocks to two. It illustrates the central distinction of this chapter: changing the storage partition is exact because every column is still represented. Dropping an operator coupling is a different operation and requires the error analysis developed next.
